% Standalone problem document, generated from the adjacent README.md.
% Compile with XeLaTeX. Mathematical content is maintained in Markdown.
\documentclass[11pt,a4paper]{article}
\usepackage[fontsize=10.5pt]{scrextend}
\usepackage[margin=22mm,headheight=15pt,headsep=9mm,footskip=11mm]{geometry}
\usepackage{fontspec}
\setmainfont{texgyrepagella}[Extension=.otf,UprightFont=*-regular,BoldFont=*-bold,ItalicFont=*-italic,BoldItalicFont=*-bolditalic]
\setsansfont{texgyreheros}[Extension=.otf,UprightFont=*-regular,BoldFont=*-bold,ItalicFont=*-italic,BoldItalicFont=*-bolditalic]
\setmonofont{texgyrecursor}[Extension=.otf,UprightFont=*-regular,BoldFont=*-bold,ItalicFont=*-italic,BoldItalicFont=*-bolditalic,Scale=MatchLowercase]
\usepackage{amsmath,amssymb}
\usepackage{unicode-math}
\setmathfont{texgyrepagella-math.otf}
\usepackage{xcolor,microtype,enumitem,needspace,fancyhdr}
\usepackage{bookmark}
\definecolor{ink}{HTML}{17344B}
\definecolor{accent}{HTML}{197D80}
\definecolor{muted}{HTML}{596773}
\hypersetup{colorlinks=true,linkcolor=accent,urlcolor=accent,pdftitle={MF-13 - Symmetric
maximizers for Lyapunov operators of order
six},pdfauthor={Open Problems in Numerical Linear Algebra}}
\urlstyle{same}
\setlength{\parindent}{0pt}
\setlength{\parskip}{5pt plus 1pt minus 1pt}
\linespread{1.04}
\setlength{\emergencystretch}{3em}
\widowpenalty=10000
\clubpenalty=10000
\displaywidowpenalty=10000
\allowdisplaybreaks[1]
\setlist{nosep,leftmargin=1.5em}
\providecommand{\tightlist}{\setlength{\itemsep}{0pt}\setlength{\parskip}{0pt}}
\providecommand{\pandocbounded}[1]{#1}
\pagestyle{fancy}
\fancyhf{}
\fancyhead[L]{\sffamily\footnotesize\color{muted}OPEN PROBLEMS IN NUMERICAL LINEAR ALGEBRA}
\fancyhead[R]{\sffamily\bfseries\footnotesize\color{accent}MF-13}
\fancyfoot[L]{\parbox[b]{0.9\textwidth}{\sffamily\fontsize{8}{10}\selectfont\color{muted}Editorial ratings; a bounded search cannot certify that a problem remains unsolved.\\Literature check: 2026-09-10}}
\fancyfoot[R]{\sffamily\footnotesize\color{muted}\thepage}
\renewcommand{\headrulewidth}{0.3pt}
\renewcommand{\headrule}{\hbox to\headwidth{\color{accent}\leaders\hrule height \headrulewidth\hfill}}
\makeatletter
\renewcommand{\section}{\Needspace{5\baselineskip}\@startsection{section}{1}{0pt}{12pt plus 2pt minus 2pt}{4pt}{\sffamily\bfseries\large\color{ink}}}
\renewcommand{\subsection}{\Needspace{4\baselineskip}\@startsection{subsection}{2}{0pt}{9pt plus 1pt minus 1pt}{3pt}{\sffamily\bfseries\normalsize\color{ink}}}
\makeatother
\setcounter{secnumdepth}{0}
\begin{document}
{\sffamily\small\color{muted}Matrix functions and stability\par}
\vspace{3pt}
{\raggedright\hyphenpenalty=10000\exhyphenpenalty=10000\sffamily\bfseries\fontsize{21}{24}\selectfont\color{ink}Symmetric
maximizers for Lyapunov operators of order six\par}
\vspace{7pt}
{\sffamily\small\textbf{Difficulty:} challenging\quad\textbf{Importance:} interesting
to specialist\par}
{\sffamily\small\bfseries\color{accent}Status: Open\par}
\vspace{4pt}
{\color{accent}\hrule height 0.8pt}
\vspace{6pt}
\textbf{Rating rationale:} Challenging because the last undecided
dimension sits between proved cases and counterexamples; specialist
impact reflects the surviving order-six norm identity rather than the
refuted all-orders claim.

\subsection{Problem statement}\label{problem-statement}

For each real matrix \(A\in\mathbb R^{6\times6}\), define
\(L_A(X)=AX+XA^T\). Write \(\|X\|_F=(\operatorname{tr}(X^TX))^{1/2}\).
Is it true that

\[
\max_{X\in\mathbb R^{6\times6},\ \|X\|_F=1}\|L_A(X)\|_F
=
\max_{X=X^T,\ \|X\|_F=1}\|L_A(X)\|_F
\qquad\text{for every }A\in\mathbb R^{6\times6}?
\]

This concerns the \textbf{largest} singular value of the Lyapunov
operator, with no stability assumption on \(A\).

\subsection{Why it matters}\label{why-it-matters}

Restricting a norm calculation to symmetric matrices substantially
reduces the associated singular-value problem. The conjecture asks when
that reduction retains the exact operator norm, an issue relevant to
matrix-equation norm estimation.

\subsection{References}\label{references}

\begin{enumerate}
\def\labelenumi{\arabic{enumi}.}
\tightlist
\item
  D. Kressner and B. Vandereycken, \emph{A counterexample to the
  symmetric-maximizer conjecture for Lyapunov operators},
  arXiv:2608.20875v1, 21 August 2026, §§1 and 5.
  \href{https://arxiv.org/html/2608.20875}{Primary text}.
\item
  R. Byers and S. Nash, \emph{On the singular ``vectors'' of the
  Lyapunov operator}, SIAM Journal on Algebraic and Discrete Methods
  8(1) (1987), 59--66. \href{https://doi.org/10.1137/0608003}{DOI}.
\end{enumerate}

\subsection{Status check --- 2026-09-08}\label{status-check-2026-09-08}

Reference 1 explicitly leaves order six open, while refuting every order
at least seven. The original all-orders assertion must therefore not be
catalogued as open. The latest arXiv record is v1. Searches included
\textit{symmetric-maximizer Lyapunov order six},
\textit{Lyapunov conjecture 2026 proof}, and the exact paper title;
no subsequent resolution of order six was located. This is a targeted
literature check, not a certification of openness.

\subsection{Audit --- 2026-09-10}\label{audit-2026-09-10}

Rechecked the \href{https://arxiv.org/html/2608.20875}{August 2026
counterexample paper, §5}: order six alone remains unresolved.
Exact-title and order-six searches found no subsequent resolution. The
importance rating now reflects this narrow surviving target; the broader
symmetric-maximizer assertion has already been refuted.
\end{document}
